% J1.1: dónde se ejecuta Julia según lo que se quiere hacer: explorar, repetir o publicar.
\documentclass[tikz,border=4pt]{standalone}
\usepackage{resonancia}
\begin{document}
\begin{tikzpicture}[node distance=2.4cm]
  \node[cajaAzul, minimum width=2.8cm] (repl) at (0,0) {\textbf{REPL}\\[-1pt]{\footnotesize probar una línea}};
  \node[cajaVerde, minimum width=2.8cm] (scr) at (4.2,0) {\textbf{Script} \texttt{.jl}\\[-1pt]{\footnotesize repetir y automatizar}};
  \node[cajaNaranja, minimum width=2.8cm] (qmd) at (8.4,0) {\textbf{Quarto} \texttt{.qmd}\\[-1pt]{\footnotesize explicar y publicar}};
  \node[caja, draw=rVioleta, fill=rVioleta!8, minimum width=2.8cm] (plu) at (4.2,-2.0) {\textbf{Pluto}\\[-1pt]{\footnotesize explorar con deslizadores}};
  \draw[flecha] (repl) -- node[etiqueta, above] {se guarda} (scr);
  \draw[flecha] (scr) -- node[etiqueta, above] {se documenta} (qmd);
  \draw[flecha, color=rGris] (repl) -- (plu);
  \draw[flecha, color=rGris] (plu) -- (scr);
  \node[etiqueta, color=rGris, anchor=east] at (-0.1,-1.0) {idea suelta};
  \node[etiqueta, color=rGris, anchor=west] at (4.3,-1.0) {se consolida};
  \node[etiqueta, anchor=north] at (4.2,-2.75) {Todos usan el mismo entorno: \texttt{julia --project=.}};
\end{tikzpicture}
\end{document}
